\subsection{The Normal Distribution}
\hrulefill

The Normal Distribution is by far the most important distribution that we will study in this class.
It is the distribution most often applied to data by scientists in all disciplines. It was first used 
in 733 by Moivre and used by Gauss to predict the location of astronomical bodies.\\

Also known as:
\begin{itemize}
    \item Normal Distribution
    \item Gaussian Distribution
    \item Bell Curve
\end{itemize}

\paragraph*{Definition}
A continuous random variable $X$ is said to have a normal distribution with mean $\mu$ and 
variance $\sigma^2$ if its PDF is given by
\[f(x) = \frac{1}{\sigma \sqrt{2 \pi}} e^{-\frac{(x - \mu)^2}{2 \sigma^2}} \quad -\infty < x < \infty.\]
We denote this as $X \sim \mathcal{N}(\mu, \sigma^2)$.

\paragraph*{Expected Value and Variance}
If $X \sim \mathcal{N}(\mu, \sigma^2)$, then
\[\mu = E(X), \quad \sigma^2 = Var(X).\]

Note that because of the symmetry of normal distributions, $\mu = \eta(0.5) = \text{median}$.

\subsubsection*{Standard Normal Distribution}
A normal random variable with mean 0 and variance 1 is called a standard normal random variable.
We denote this as $Z \sim \mathcal{N}(0, 1)$ with CDF $\Phi(z)$.

\paragraph*{Remarks}
\begin{itemize}
    \item The standard normal distribution does not often arise naturally, but it serves as a tool for computing probabilities involving the normal distribution.
    \item Tables of values of $\Phi(z)$ are used to find probabilities for normal random variables.
\end{itemize}

\paragraph*{Standard Normal Table}
A standard normal table gives values of $\Phi(z)$ for various values of $z$.
To use the table, we find the row corresponding to the first two digits of $z$ and the column corresponding to the second decimal place of $z$.
The value at the intersection gives $\Phi(z)$.

\paragraph*{Normal Percentiles}
The $k$th percentile $\eta(k)$ of a normal distribution is the value $z_k$ such that $\Phi(z_k) = k$.
To find $z_k$, we look up the value $k$ in the body of the standard normal table and find the corresponding $z$ value from the row and column headers.

\paragraph*{Tail Probabilities}
Later in this course we will frequently work with right tail areas of the standard normal distribution. Let 
$z_a$ denote the value on the measurement axis, for which the area under the standard normal curve to the right of $z_a$
is equal to $a$.
\[
    z_a = (1-a)100 \text{-th percentile}
\]

\paragraph*{Standardization}
\textbf{Any} normal random variable $X$ with mean $\mu$ and standard deviation $\sigma$ can be transformed into a random variable $Z$ 
having a standard normal distribution via the standardization transformation:
\[
    Z = \frac{X - \mu}{\sigma}.
\]

This allows us to compute probabilities for any normal random variable using the standard normal table.
Let $X \sim \mathcal{N}(\mu, \sigma^2)$, then for any $a$ and $b$,
\[P(a \leq X \leq b) = P\left(\frac{a - \mu}{\sigma} \leq Z \leq \frac{b - \mu}{\sigma}\right) = \Phi\left(\frac{b - \mu}{\sigma}\right) - \Phi\left(\frac{a - \mu}{\sigma}\right).\]

\pagebreak

\paragraph*{Example}
Let $X=$ height of a random female student entering college in inches, where $X \sim \mathcal{N}(\mu=65, \sigma^2=25)$.\\
\\
a. What is the probability that a randomly selected female student has a height between 55 and 72.5 inches?
\begin{align*}
    P(55 \leq X \leq 72.5) &= P\left(\frac{55 - 65}{5} \leq Z \leq \frac{72.5 - 65}{5}\right) \\
    &= P(-2 \leq Z \leq 1.5) \\
    &= \Phi(1.5) - \Phi(-2) \\
    &= 0.9332 - 0.0228 \\
    &= 0.9104.
\end{align*}

b. What is the height of a student such that only 5\% of the students are taller than her? (Find the 95th percentile.)
\begin{align*}
    P(X > h) &= 0.05 \\
    P\left(Z > \frac{h - 65}{5}\right) &= 0.05 \\
    P\left(Z < \frac{h - 65}{5}\right) &= 0.95 \\
    \frac{h - 65}{5} &= \Phi^{-1}(0.95) \\
    \frac{h - 65}{5} &= 1.645 \\
    h - 65 &= 8.225 \\
    h &= 73.225.
\end{align*}

\pagebreak

\paragraph*{Example:}
LSAT scores have an approximately Normal distribution with mean $\mu = 150$ and standard deviation $\sigma = 10$.\\
\\
a. If you have a score of 167, then you did better than what percent of the people taking the test with you?
\begin{align*}
    P(X < 167) &= P\left(Z < \frac{167 - 150}{10}\right) \\
    &= P(Z < 1.7) \\
    &= \Phi(1.7) \\
    &\approx 0.9554.
\end{align*}

b. How high would your score have to be so that you are in the top 90\%?
\begin{align*}
    P(X > s) &= 0.10 \\
    P\left(Z > \frac{s - 150}{10}\right) &= 0.10 \\
    P\left(Z < \frac{s - 150}{10}\right) &= 0.90 \\
    \frac{s - 150}{10} &= \Phi^{-1}(0.90) \\
    \frac{s - 150}{10} &\approx 1.28 \\
    s - 150 &\approx 12.8 \\
    s &\approx 162.8.J
\end{align*}